\subsection{Finite dimensionality of the space of nondegenerate short star-products}

\subsubsection{The case of a semisimple automorphism}

\begin{Proposition}\label{finde1} Let $A$ be a finitely generated Poisson algebra such that the Poisson scheme $X=\operatorname{Spec}A$ has finitely many symplectic leaves. Let~${\mathbf A}$ be a quantization of~$A$, and \mbox{$g\colon {{\mathbf A}}\to {{\mathbf A}}$} be a semisimple automorphism which commutes with~$s$. Then nondegenerate short star-products on~$A$, if they exist, are parametrized by points of the finite dimensional vector space \linebreak $(HH_0({{\mathbf A}},{{\mathbf A}}g)^s)^*$ which don't belong to a countable collection of hypersurfaces, modulo scaling.

Moreover, $\dim HH_0({{\mathbf A}},{{\mathbf A}}g)^s\le \dim HP_0(X^g)^s<\infty$, where $X^g$ is the fixed point scheme of~$g$ on~$X$, and $HP_0(X^g):= \mathcal{O}(X^g)/\lbrace{\mathcal{O}(X^g), \mathcal{O}(X^g)\rbrace}$ is the zeroth Poisson homology of $X^g$.\footnote{Here, abusing notation, we denote the associated graded isomorphism $\operatorname{gr}g\colon A\to A$ simply by $g$.}
\end{Proposition}

\begin{proof} Let $E$ be the span of elements ${\mathbf{ab}}-{\mathbf{b}}g({\mathbf a})$ for ${\mathbf{a}},{\mathbf{b}}\in {{\mathbf A}}$.

We claim that for any $a,b\in A$, $\operatorname{gr}(E)$ contains the element $(a-g(a))b$. Indeed, we may assume that $a$, $b$ are homogeneous of degrees $i$, $j$, and let ${{\mathbf a}}$, ${\mathbf b}$ be their lifts to ${{\mathbf A}}$. Then the leading term (of degree $i+j$) of ${{\mathbf a}} {\mathbf b}-{\mathbf b}g({{\mathbf a}})$ is $(a-g(a))b$, as desired.

\looseness=-1 Also, we claim that if $a\in A$ and $g(a)=a$ then for any $b\in A$, $\lbrace{a,b\rbrace}$ belongs to $\operatorname{gr}(E)$. Indeed, we can again assume that $a$, $b$ are homogeneous of degrees $i$, $j$. Since $g=1$ on its generalized 1-eigenspace, we can choose a lift ${\mathbf a}$ of~$a$ to~${\mathbf A}$ which is invariant under~$g$. Then choosing any lift~${\mathbf b}$ of~$b$ to~${\mathbf A}$, we see that the leading term (of degree $i+j-2$) of ${{\mathbf a}} {\mathbf b}-{\mathbf b}g({{\mathbf a}})=[{{\mathbf a}},{\mathbf b}]$ is $\lbrace{ a,b\rbrace}$, as desired.

Now let $I\subset A$ be the span of elements $(a-g(a))b$. This is the ideal of the fixed point subscheme~$X^g$. Note that $I\cap A^g$ is a Poisson ideal in $A^g$. Hence, $A/I=A^g/(I\cap A^g)$ is a Poisson algebra, so~$X^g$ is a Poisson scheme. Moreover, since $g=1$ on its generalized 1-eigenspace, the $g$-invariants in the tangent space to the symplectic leaf~$L$ at any point $P\in X^g\cap L$ form a~nondegenerate subspace, hence $X^g\cap L$ is symplectic. Thus, $X^g$ also has finitely many symplectic leaves, which are connected components of intersections of $X^g\cap L$ for various~$L$.

 Now, the space $A/\operatorname{gr}(E)$ receives a surjective $s$-invariant map from the space $HP_0(X^g)=\mathcal{O}(X^g)/\lbrace{ \mathcal{O}(X^g),\mathcal{O}(X^g) \rbrace}$. This space is finite dimensional by \cite[Theorem~1.4]{ES}. Hence the space $A/\operatorname{gr}(E)$
 is finite dimensional. Therefore the space ${{\mathbf A}}/E=HH_0({{\mathbf A}},{{\mathbf A}}g)$ is finite dimensional, and we have the inequalities
\begin{gather*}
\dim HH_0({{\mathbf A}},{{\mathbf A}}g)\le \dim HP_0(X^g),\qquad \dim HH_0({{\mathbf A}},{{\mathbf A}}g)^s\le \dim HP_0(X^g)^s.
\end{gather*}

\looseness=-2 Now the proposition follows, since by Corollary~\ref{bun}, short nondegenerate star-products corre\-spon\-ding to~$g$ are parametrized by the subset of $(HH_0({{\mathbf A}},{{\mathbf A}}g)^s)^*$ determined by the nondegenera\-cy condition, which is the vanishing condition for countably many determinant polynomials.
\end{proof}

\begin{Remark} \quad
\begin{enumerate}\itemsep=0pt
\item It is clear from the proof of Proposition~\ref{finde1} that it holds more generally, namely when~$g$ acts trivially on its generalized 1-eigenspace (this is the only consequence of semisimplicity of~$g$ used in the proof).
\item In particular, Proposition \ref{finde1} applies to the setting of \cite{L2}, namely $A=\mathcal{O}(X)$, where~$X$ is a conical symplectic singularity in the sense of Beauville. Indeed, in this case it is known that~$X$ has finitely many symplectic leaves \cite[Theorem~2.3]{Ka}. This property also holds for Higgs branches considered in~\cite{BPR} (by \cite[Section~2.3]{L3}), and is expected for Coulomb branches which are cones, so Proposition~\ref{finde1} applies to all such cases.
\end{enumerate}
\end{Remark}
